\section{Conclusión}

Los modelos de Inteligencia Artificial analizados muestran distintas formas de resolver problemas relacionados con percepción, aprendizaje, razonamiento, clasificación, predicción, planificación y toma de decisiones. Los agentes inteligentes permiten comprender la IA como un proceso de interacción entre un sistema y su ambiente. Los agentes de reflejo simple son útiles para tareas rápidas y controladas, mientras que los agentes basados en metas y utilidad permiten resolver problemas más complejos mediante planificación y comparación de alternativas.

La representación de problemas y los algoritmos de búsqueda muestran que muchas situaciones pueden formularse como espacios de estados. Los sistemas expertos destacan por su capacidad de explicación, aunque son rígidos si el conocimiento cambia. Los modelos probabilísticos permiten actuar en condiciones de incertidumbre. El aprendizaje automático y las redes neuronales hacen posible descubrir patrones en grandes volúmenes de datos. El aprendizaje por refuerzo permite optimizar decisiones secuenciales. Finalmente, la visión por computadora, el procesamiento del lenguaje natural y la robótica muestran cómo la IA se aplica en sistemas reales que perciben, interpretan, comunican y actúan.

Una de las principales dificultades encontradas al estudiar estos modelos es comprender que cada técnica tiene condiciones específicas de uso. Algunos modelos requieren muchos datos; otros necesitan reglas bien definidas; otros dependen de buenas heurísticas; y otros sacrifican interpretabilidad a cambio de mayor desempeño. Por ello, seleccionar un modelo de IA no consiste únicamente en elegir el algoritmo más avanzado, sino en analizar el problema, el ambiente, los recursos disponibles, los riesgos y los criterios de evaluación.

Como aprendizaje principal, se concluye que la Inteligencia Artificial es un campo integrador. Un sistema inteligente de calidad combina percepción, representación, razonamiento, aprendizaje y acción. Además, su desarrollo debe considerar no solo la precisión técnica, sino también la seguridad, la transparencia, la ética y el impacto social. En este sentido, estudiar los modelos de IA permite no solo construir sistemas más capaces, sino también evaluar críticamente sus alcances y limitaciones en problemas reales.

